\chapterzero

\begin{learningobjectives}
By the end of this chapter, you should be able to:
\begin{itemize}
    \item translate ordinary mathematical statements into propositional and predicate logic;
    \item construct truth tables and test logical equivalence;
    \item distinguish implication, converse, inverse, and contrapositive;
    \item read quantified statements and write their negations correctly;
    \item recognize the logical patterns used in direct proof, contradiction, and induction.
\end{itemize}
\end{learningobjectives}

\begin{mynote}
This chapter is optional. Readers already comfortable with logical connectives, quantifiers, and standard proof patterns may proceed directly to Chapter 2 and return here as needed.
\end{mynote}

Mathematics becomes reliable only when its statements are made precise. This chapter assembles the language needed for that precision: propositional logic, predicate logic, and the standard patterns of mathematical proof. The material is elementary, but it is used constantly — every definition, hypothesis, and proof later in these notes is written in this language.

A note on scope: this chapter treats logic as a working tool rather than a subject in itself. We state what a proposition, a quantifier, and a valid inference are, and we practice translating ordinary mathematical sentences into symbols. Questions about the axiomatization of logic itself belong to Chapter 1.

\section{Propositional Logic}

\subsection{Naïve introduction to propositional logic}
Before we start the main part of the propositional logic, let’s start with a question: what is a proposition? Take a look at the sentences below:
\begin{enumerate}
    \item The equation $x^2+1=0$ has a real-number solution;
    \item There are infinitely many prime numbers;
    \item Every even integer $n \ge 4$ is a sum of two prime numbers (The famous Goldbach Conjecture);
    \item What time is it?
    \item This statement is false;
\end{enumerate}
In the sentences above, the first 3 sentences share the same properties: 1) they are sentences or assertions that declare facts. 2) They are either true or false, but not both. However, the 4th sentence is not a declarative sentence, and we can’t judge the correctness of the 5th sentence.

Mathematics is built from sentences that assert facts and are therefore true or false. Whether a given statement is actually true is a separate question; what matters first is that it has a determinate truth value. Other kinds of sentences have their uses, but axioms, definitions, and theorems are all written in this declarative form, so it is the form we must learn to manipulate.

\begin{definition}[Proposition]
Propositions are sentences or assertions that declare facts and are either true or false, but not both.
\end{definition}

To make propositions more actionable, we use \textbf{propositional variables} to represent different propositions. Commonly used propositional variables are letters like $p, q, r, \dots$

Sometimes we need to connect different propositional variables to form a new proposition. And we need to use words like “and”, “or”, “imply” to explain their relationships.

\begin{definition}[Logical operators/connectives]
Logical operators/connectives are marks that are used to connect propositional variables, forming compound propositions.
\end{definition}

Here are some common operators:
\begin{description}
    \item[$\neg$] negation (“not”)
    \item[$\wedge$] conjunction (“and”)
    \item[$\vee$] disjunction (“or”)
    \item[$\to$ (or $\Rightarrow$)] conditional (“imply”)
    \item[$\leftrightarrow$] biconditional (“if and only if” (“iff”))
    \item[$\oplus$] exclusive-or
\end{description}

\begin{definition}[Compound proposition]
A compound proposition is built from propositional variables or constants through logical operators.
\end{definition}

\begin{definition}[Atomic proposition]
An atomic proposition is a proposition that can't be divided any further. It is the most basic building block of logical expressions and is considered an indivisible semantic unit.
\end{definition}

Whether an atomic proposition is true depends on subject matter outside logic. The truth of a compound proposition, however, is determined by the truth of its parts alone, and truth tables make this dependence explicit.

\begin{definition}[Truth table]
A truth table is a mathematical table used in logic to compute the functional values of logical expressions based on their inputs.
\end{definition}

\paragraph{Negation}
\begin{center}
\begin{tabular}{c|c}
\toprule
$p$ & $\neg p$ \\
\midrule
T & F \\
F & T \\
\bottomrule
\end{tabular}
\end{center}

\paragraph{Conjunction}
\begin{center}
\begin{tabular}{cc|c}
\toprule
$p$ & $q$ & $p \wedge q$ \\
\midrule
T & T & T \\
T & F & F \\
F & T & F \\
F & F & F \\
\bottomrule
\end{tabular}
\end{center}



\paragraph{Disjunction}
\begin{center}
\begin{tabular}{cc|c}
\toprule
$p$ & $q$ & $p \vee q$ \\
\midrule
T & T & T \\
T & F & T \\
F & T & T \\
F & F & F \\
\bottomrule
\end{tabular}
\end{center}

\paragraph{Conditional (Imply)}
\begin{center}
\begin{tabular}{cc|c}
\toprule
$p$ & $q$ & $p \to q$ \\
\midrule
T & T & T \\
T & F & F \\
F & T & T \\
F & F & T \\
\bottomrule
\end{tabular}
\end{center}

\paragraph{Biconditional (iff)}
\begin{center}
\begin{tabular}{cc|c}
\toprule
$p$ & $q$ & $p \leftrightarrow q$ \\
\midrule
T & T & T \\
T & F & F \\
F & T & F \\
F & F & T \\
\bottomrule
\end{tabular}
\end{center}

\paragraph{Exclusive-or}
\begin{center}
\begin{tabular}{cc|c}
\toprule
$p$ & $q$ & $p \oplus q$ \\
\midrule
T & T & F \\
T & F & T \\
F & T & T \\
F & F & F \\
\bottomrule
\end{tabular}
\end{center}

Just like numerical computation, logical operators also have precedence. Here is the order of precedence for logical operators (from highest to lowest):
$$ ( ) \quad [ ] \quad \neg \quad \wedge \quad \vee \quad \oplus \quad \to \quad \leftrightarrow $$
When operators of the same precedence appear together, association is taken from left to right, and parentheses always override this table.

Using these laws, we can define the truth value of different propositions. But there exist some sub-classes of compound propositions:
\begin{description}
    \item[Tautology] A tautology is a compound proposition that is always true, no matter what the truth values of the propositional variables are.
    \item[Contradiction] A compound proposition that is always false, regardless of the truth values of propositional variables.
    \item[Contingency] A compound proposition that is neither a tautology nor a contradiction. It can be true or false, depending on the value of the variables.
\end{description}

Based on the definitions above, now we can rigorously define what logical equivalence is.

\begin{definition}[Logical Equivalence]
We can say compound propositions $p, q$ are \textbf{logically equivalent} if $p \leftrightarrow q$ is a tautology. If $p$ and $q$ are logically equivalent, we denote it as $p \equiv q$.
\end{definition}

In other words, $p$ and $q$ receive the same truth value under every truth assignment of their atomic propositions. Section 0.2 shows that a small set of connectives is already enough to express every compound proposition.

The following laws are used constantly, and they are worth internalizing rather than looking up each time.
\begin{align*}
    \text{Identity Laws:} \quad & p \wedge T \equiv p \\
    & p \vee F \equiv p \\
    \text{Domination Laws:} \quad & p \vee T \equiv T \\
    & p \wedge F \equiv F \\
    \text{Absorption Laws:} \quad & p \vee (p \wedge q) \equiv p \\
    & p \wedge (p \vee q) \equiv p \\
    \text{Negation Laws:} \quad & p \vee \neg p \equiv T \\
    & p \wedge \neg p \equiv F \\
    \text{Commutative Laws:} \quad & p \vee q \equiv q \vee p \\
    & p \wedge q \equiv q \wedge p \\
    \text{Associative Laws:} \quad & (p \vee q) \vee r \equiv p \vee (q \vee r) \\
    & (p \wedge q) \wedge r \equiv p \wedge (q \wedge r) \\
    \text{Distributive Laws:} \quad & p \vee (q \wedge r) \equiv (p \vee q) \wedge (p \vee r) \\
    & p \wedge (q \vee r) \equiv (p \wedge q) \vee (p \wedge r) \\
    \text{De Morgan’s Laws:} \quad & \neg(p \wedge q) \equiv \neg p \vee \neg q \\
    & \neg(p \vee q) \equiv \neg p \wedge \neg q
\end{align*}

\begin{definition}[Satisfiability]
A compound proposition is \textbf{satisfiable} if it is true under some truth assignment for propositional variables. A truth assignment that makes a compound proposition true is called a \textbf{solution}. If a compound proposition is not satisfiable, we say it is \textbf{unsatisfiable}.
\end{definition}

For small formulas, a truth table decides satisfiability directly. For large ones this is impractical, and deciding satisfiability becomes the Boolean satisfiability problem (SAT), a central problem of computer science for which specialized algorithms exist.

\subsection{Functional Completeness}
The connectives introduced above form compound propositions, but nothing so far suggests that all of them are necessary. In fact a surprisingly small collection suffices: two connectives can already express every compound proposition. This section makes that claim precise through normal forms.

\begin{definition}[Disjunctive Normal Form]
A compound proposition is in \textbf{disjunctive normal form (DNF)} if it is a disjunction of conjunctions of propositional variables or their negations.
\end{definition}

\begin{theorem}\label{thm:dnf}
Every compound proposition is logically equivalent to some compound proposition in disjunctive normal form.
\end{theorem}
\begin{proof}
The proof is constructive: given any compound proposition, write its truth table, then read off the rows in which it is true. Each such row contributes one conjunction, and the disjunction of these conjunctions is the proposition in disjunctive normal form. The truth table of the exclusive-or connective illustrates the construction:
\begin{center}
\begin{tabular}{cc|c}
\toprule
$p$ & $q$ & $p \oplus q$ \\
\midrule
T & T & F \\
T & F & T \\
F & T & T \\
F & F & F \\
\bottomrule
\end{tabular}
\end{center}
The content of the middle two lines (where the result is T) represents the solution to this compound proposition. Then we can write a proposition like this:
$$ p \oplus q \equiv (p \wedge \neg q) \vee (\neg p \wedge q) $$
Each conjunction in the disjunction records one row of the truth table, so the construction applies to every compound proposition, which proves the theorem.
\end{proof}

A dual construction, built from the rows in which the proposition is false, yields a conjunctive normal form.

\begin{definition}[Conjunctive Normal Form]
A compound proposition is in \textbf{conjunctive normal form (CNF)} if it is a conjunction of disjunctions of propositional variables or their negations.
\end{definition}

\begin{theorem}
Every compound proposition is logically equivalent to some compound proposition in conjunctive normal form.
\end{theorem}
\begin{proof}
Let $A$ be a compound proposition. By Theorem~\ref{thm:dnf}, $\neg A$ has a disjunctive normal form. Negating that expression and applying De Morgan’s laws pushes every negation onto the variables, leaving a conjunction of disjunctions equivalent to $A$.
\end{proof}

These two theorems have practical weight. A circuit designer cannot build a distinct component for every connective, and the theorems say that none are needed beyond a small stock: the collection $\{\neg, \vee, \wedge\}$ already constructs a proposition logically equivalent to any compound proposition whatsoever. This observation motivates the following definition.

\begin{definition}[Functionally Complete]
A collection of logical operators is \textbf{functionally complete} if any compound proposition is logically equivalent to a compound proposition that involves only the logical operators in the collection.
\end{definition}

Thus, we can declare that the collection $\{\neg, \vee, \wedge\}$ is functionally complete. Also, according to De Morgan’s Law:
\begin{align*}
    p \wedge q \equiv \neg(\neg p \vee \neg q) \\
    p \vee q \equiv \neg(\neg p \wedge \neg q)
\end{align*}
Conjunction is expressible through disjunction and negation, and disjunction through conjunction and negation, so both collections $\{\neg, \vee\}$ and $\{\neg, \wedge\}$ are functionally complete as well.

\begin{remark}
The extreme case is a single connective that is functionally complete by itself: negation combined with either NAND or NOR suffices. We will not use such connectives later, but the fact shows how far the reduction can go.
\end{remark}

\section{Predicate Logic}
Propositional logic treats each statement as an indivisible unit, and this is also its limitation. The statement “every even integer greater than $2$ is a sum of two primes” is a single proposition, but its content is about every even integer, and propositional logic cannot look inside it. Predicate logic extends the language so that such statements can be expressed and reasoned about.

\begin{definition}[Predicate]
A \textbf{predicate} is in the form of a statement with variables. A predicate is also known as a \textbf{propositional function}. A predicate can be denoted as $P(x, y, z, \dots)$, and $x, y, z$ are known as the variables.
\end{definition}

\begin{definition}[Quantifiers]
We use \textbf{quantifiers} to indicate the quantity of elements being referred to. We have two types of quantifiers, $\forall$ and $\exists$.
\begin{description}
    \item[$\forall$] The \textbf{universal quantifier} $\forall$ means “for all”.
    \item[$\exists$] The \textbf{existential quantifier} $\exists$ means “there exists”.
\end{description}
\end{definition}

For example, the statement “Every SJTUer is a talent” can be written in the form of predicate logic: $\forall x (\text{SJTUer}(x) \to \text{Talent}(x))$.

The variables of a predicate range over a specified set, called the \textbf{domain} of the variables. Once the domain is fixed, a predicate becomes a proposition for each concrete value, and over a finite domain $\{x_1, x_2, \dots, x_n\}$ the quantified forms reduce to:
$$ \forall x P(x) \equiv P(x_1) \wedge P(x_2) \wedge \dots \wedge P(x_n) $$
$$ \exists x P(x) \equiv P(x_1) \vee P(x_2) \vee \dots \vee P(x_n) $$
Over an infinite domain no such finite expansion exists, and the quantifier notation becomes essential.
\begin{itemize}
    \item $\forall x P(x)$ stands for “For every $x$ in domain, $P(x)$”. If there is a value $a$ in the domain such that $P(a)$ is false, the statement is false. $a$ is called a \textbf{counterexample}.
    \item $\exists x P(x)$ stands for “There exists an $x$, $P(x)$”. If for all values $a$ in the domain, $P(a)$ is false, the statement is false.
\end{itemize}

A formula counts as a statement only when every variable is accounted for. In $\forall x (P(x) \wedge Q(y))$, the variable $x$ is governed by a quantifier, but the truth value of $Q(y)$ varies with $y$, so the whole formula has no determinate truth value.

The distinction is this: a variable in the scope of a quantifier is \textbf{bound}; a variable that is not, like $y$ above, is \textbf{free}. A formula is a statement precisely when it contains no free variables.

Predicate formulas can be combined with the same connectives as propositional ones, and the quantifiers obey laws of their own.

\paragraph{De Morgan's Laws for Quantifiers}
\begin{align*}
    \neg \forall x P(x) \equiv \exists x \neg P(x) \\
    \neg \exists x P(x) \equiv \forall x \neg P(x)
\end{align*}

\paragraph{Commutative Laws for Like Quantifiers}
\begin{align*}
    \forall x \forall y P(x,y) \equiv \forall y \forall x P(x,y) \\
    \exists x \exists y P(x,y) \equiv \exists y \exists x P(x,y)
\end{align*}



\begin{note}
Like quantifiers commute; unlike quantifiers do not. Over $\mathbb{N}$, take $P(x,y)$ to mean $y > x$: the statement $\forall x \exists y P(x,y)$ is true, since every natural number has a larger one, while $\exists y \forall x P(x,y)$ is false, since no single $y$ exceeds every $x$.
\end{note}

\begin{definition}[Logical Equivalence in Predicate Logic]
If $\phi$ and $\psi$ are statements with predicates and quantifiers, but without free variables, then we say that $\phi$ and $\psi$ are \textbf{logically equivalent}, written as $\phi \equiv \psi$, if, no matter which concrete predicates (for the predicate symbols) and domains (for variables) are given, the truth values of $\phi, \psi$ coincide.
\end{definition}

Predicate logic also has its analogues of tautology (validity under every interpretation) and satisfiability, but a systematic treatment belongs to a course on mathematical logic, and we stop at the definitions.

The chapter has assembled a working language: propositions and connectives to compose statements, predicates and quantifiers to make claims about collections, and standard laws for manipulating both. Later chapters use this language constantly — most visibly in the definitions and hypotheses of Chapter 1, and quietly in every proof thereafter.

\vspace{1cm}
\noindent
\textbf{Keywords:} propositional logic, predicate logic, logical equivalence \\
\textbf{Reference:} Discrete Mathematics and Its Applications (Eighth Edition), Kenneth H. Rosen, McGraw-Hill Education.

\cleardoublepage